% SGA 3，第 VII 篇_B，英文重建本——第 1 轮（中文译本）。
% 范围：依据原文件第 3--4 页，从 0.2 至推论 0.2.D；止于该推论的证明之前。
% 合并阅读版页码： 516--517. 再版印刷页码： 505--506.
% 本范围内的原版页边定位： 480.
% 依据： Polo--Gille Exp7B-13oct24.pdf,
% SHA-256 D303690E0195654BBE7C3DAA5D89EA23E24CD709A3575217234AC19AE8836694.
% 曾参阅 Jacob Reinhold 的 CC BY 4.0 Markdown，作为对比版本来源。

\medskip
\noindent\textbf{0.2.}
\phantomsection
\label{sga3:VIIB:0.2}
设 \(A\) 是交换伪紧环。一个\emph{伪紧 \(A\)-模} \(M\) 是分离完备的
拓扑 \(A\)-模，它具有由满足 \(M/M'\) 在 \(A\) 上长度有限的子模
\(M'\) 组成的 \(0\) 的邻域基。

若 \(M\) 和 \(N\) 是伪紧 \(A\)-模，则从 \(M\) 到 \(N\) 的\emph{态射}，
按定义是连续的 \(A\)-线性映射。我们用
\[
  \operatorname{Hom}_c(M,N)
\]
表示这些态射组成的群。

\medskip
\noindent\SGAInlineRefTarget{sga3:vii:proposition:0-2-b}\textbf{命题 0.2.B。}\footnote{编者将本段的结果单独列为以下命题，然后说明证明的步骤；参见\ \textup{[CA]}，
\S~IV.3，或 \textup{[DG70]}，\S~V.2。}
\phantomsection
\label{sga3:VIIB:0.2.B}
{\itshape
\begin{enumerate}
  \item[\textup{(i)}]
  伪紧 \(A\)-模构成阿贝尔范畴，记作 \(\operatorname{PC}(A)\)。
  特别地，对每个态射 \(f:M\to N\)，子模 \(\operatorname{Im}(f)\) 都是完备的，
  从而在 \(N\) 中是闭的。

  \item[\textup{(ii)}]
  有限长度的伪紧 \(A\)-模（它们的拓扑因而是离散的）构成
  \(\operatorname{PC}(A)\) 的一族余生成元。

  \item[\textup{\SGAInlineRefTarget{sga3:vii:proposition:0-2-b:item:iii}(iii)}]
  无穷积和滤过投射极限都是正合的；换言之，\(\operatorname{PC}(A)\)
  满足公理 \((\mathrm{AB}5^{*})\)。\footnote{参见\ \textup{[Gr57]}，I，\S~\SGARefLink{sga3:VIIA:1.5}{1.5}
  和命题~1.8。例如也可参阅 \textup{[Po73]}，\S~\SGARefLink{sga3:VIIB:2.8}{2.8}，
  或 \textup{[We95]}，附录~A.4。}
\end{enumerate}}

为方便读者，我们简述证明的步骤。首先使用下列引理
（\textup{[CA]}，\S~IV.3，引理~1；证明见 \textup{[BEns]}，III，
\S~\SGARefLink{sga3:VIIA:7.4}{7.4}，定理~1 和例~2）。

\medskip
\noindent\textbf{引理 0.2.C。}
\phantomsection
\label{sga3:VIIB:0.2.C}
\textit{设 \(B\) 是环，\(I\) 是滤过有序集，\((M_i)\) 和 \((N_i)\)
是两个以 \(I\) 为指标的左 \(B\)-模投射系。设
\[
  (s_i):(M_i)\longrightarrow(N_i)
\]
是投射系的态射，且对每个 \(i\)，\(s_i\) 都是满射并具有 Artin 核。那么映射}
\[
  \varprojlim_i s_i:
  \varprojlim_i M_i
  \longrightarrow
  \varprojlim_i N_i
\]
\textit{是满射。}

\medskip
\noindent\textbf{推论 0.2.D}
\textit{(\textup{[CA]}, \S~IV.3, 命题~10 和~11).}
\phantomsection
\label{sga3:VIIB:0.2.D}
{\itshape 设 \(M\) 是伪紧 \(A\)-模。
\begin{enumerate}
  \item[\textup{(i)}]
  若 \(K\) 是 \(M\) 的闭子模，则赋予商拓扑的 \(M/K\) 是伪紧 \(A\)-模。

  \item[\textup{(ii)}]
  设 \((M_i)\) 是 \(M\) 的闭子模的滤过递减族。
  \begin{enumerate}
    \item[\textup{\SGAInlineRefTarget{sga3:VIIB:0.2.D:item:a}(a)}]
    典范映射
    \[
      M\longrightarrow\varprojlim_i M/M_i
    \]
    是满射，且其核为 \(\bigcap_i M_i\)。

    \item[\textup{\SGAInlineRefTarget{sga3:VIIB:0.2.D:item:b}(b)}]
    对于 \(M\) 的每个闭子模 \(N\)，有
    \[
      N+\bigcap_i M_i
      =
      \bigcap_i(N+M_i).
    \]
  \end{enumerate}
\end{enumerate}}

% 下一断点： Exp7B 原文件页码 4, 合并阅读版页码 517,
% 再版印刷页码 506, 推论 0.2.D 的证明处。
